\section{The checked algebraic backend}\label{sec:algebra-backend}

\subsection{Why this backend matters}

Many packets reduce a theorem to a finite polynomial identity but then stop at a semantic \texttt{ring} proof.  The current \RepoPath{Derivations} library closes this gap for a substantial class of identities.  It derives arithmetic laws in the embedded theory, quotients terms by provable equality, embeds the quotient semiring into a ring of differences, performs metalevel ring normalization, and reflects the normalized equality back to \(\Prf\,\PRAAx\).

This architecture is exactly compatible with \KIFCR.  An algebra node stores an output equation and a proof-producing tactic invocation.  The equation can then be consumed by later search, loop, or observer nodes through ordinary equality substitution.

\subsection{Provable-equality quotient}

Fix a context \(\Gamma\).  Define \(s\sim_\Gamma t\) when \(\Gamma\vdash s=t\).  Reflexivity, symmetry, and transitivity are derived in the deduction system; congruence follows from equality substitution.  Thus \(\sim_\Gamma\) is a congruence for the derived addition and multiplication symbols.  The quotient \(Q_\Gamma\) is a cancellative commutative semiring because the corresponding laws and cancellation are proved in \iPRA.

Natural-number terms embed into the Grothendieck ring \(G(Q_\Gamma)\).  Injectivity of the canonical map follows from additive cancellation.  An integer term is represented by a pair \((p,n)\) denoting \(p-n\), and equality of integer terms is translated to the natural equality
\[
 p+n'=p'+n.
\]
Consequently, if a polynomial identity holds in \(G(Q_\Gamma)\), injectivity yields an equality in \(Q_\Gamma\), which by definition is an object-level derivation.

\subsection{Sound use of metalevel tactics}

The tactic \texttt{ipra\_ring} does not assert that ring normalization is an axiom of \iPRA.  Its proof term has the following shape:
\begin{enumerate}[label=(\arabic*)]
  \item map each object term to its class in \(Q_\Gamma\) and then to \(G(Q_\Gamma)\);
  \item use Lean's ring tactic to prove equality in this metalevel ring;
  \item apply injectivity to recover equality of quotient classes;
  \item unfold quotient equality to obtain \(\Gamma\vdash s=t\).
\end{enumerate}
The current axiom-audit module reports no \texttt{sorryAx}.  It lists only Lean's standard metatheoretic axioms for quotients, choice, and propositional extensionality; the object theorem still depends only on \(\PRAAx\).

\subsection{Algebra node interface}

A generic algebra node can be represented by:
\[
  \mathsf{AlgNode}(\Gamma,s,t,h),
  \qquad h:\Gamma\vdash[\PRAAx]s=t.
\]
Its operational output may be one side of the equation; its proof output is \(h\).  Common subclasses are:
\begin{itemize}
  \item polynomial normalization and factorization;
  \item linear-combination certificates for ideal membership;
  \item determinant and adjugate identities at fixed dimension;
  \item integer-coordinate geometry identities;
  \item finite combinatorial sum identities after recurrence expansion.
\end{itemize}

Not every algebraic theorem is a ring identity.  Inequalities require order lemmas and often sums of squares or explicit witnesses.  Divisibility requires a witness term.  Matrix statements of variable dimension require loop invariants.  The node interface is nevertheless useful: these stronger nodes may call the ring backend for their local algebraic subgoals.

\subsection{Worked inherited examples}

The current repository supplies actual embedded derivations for:
\begin{enumerate}[label=(\roman*)]
  \item \((xz+yw)^2+(xw-yz)^2=(x^2+y^2)(z^2+w^2)\);
  \item the Heron factorization
  \[
  4b^2c^2-(b^2+c^2-a^2)^2
   =(a+b+c)(-a+b+c)(a-b+c)(a+b-c);
  \]
  \item the coordinate identity equivalent to Heron's formula;
  \item the coordinate law of cosines.
\end{enumerate}
The accompanying Lean adapter wraps the open two-square and Heron derivations as \KIFCR\ kernel rules.  This is an inherited checked component, not a new unverified claim of this packet.

\subsection{Extending the backend}

A practical extension order is:
\begin{enumerate}[label=(\arabic*)]
  \item derive a reusable ordered-semiring layer and a tactic for polynomial inequalities from explicit nonnegative decompositions;
  \item add finite-sum normalization, with a loop invariant connecting recursive sums to closed forms;
  \item add fixed-size matrix identities by flattening entries and calling \texttt{ipra\_ring};
  \item add variable-size matrix algorithms as certified loops with determinant/minor certificates;
  \item add rational-code fields under explicit nonzero denominator guards.
\end{enumerate}
Each extension should expose object-level theorem constructors rather than a semantic API only.
